\section{Smoothing a bounded positive potential}
\label{app:smoothing}

We prove Lemma~\ref{nef:bounded-smoothing}. The problem is to smooth
a bounded potential while losing arbitrarily little positivity.
Bounded plurisubharmonic functions need not be continuous, so uniform
convergence of their regularizations is unavailable. Instead, we
compare regularizations at fixed multiples of the same small radius.
The comparison becomes uniformly small and permits gluing by a
regularized maximum.

\begin{lemma}[Fixed-factor changes of radius]
\label{nef:radius-comparison}
Let $\psi$ be plurisubharmonic on a domain in $\C^n$.  Suppose that
all closed balls $\overline B(x,R)$ with centers in a fixed set $E$
are contained in the domain, and that $m\leq\psi\leq M$ on their
union.  Put
\[
 S_x(r)=\sup_{|z-x|\leq r}\psi(z),\qquad 0<r\leq R.
\]
For every $A>1$ and $Ar<R$,
\begin{equation}
 0\leq S_x(Ar)-S_x(r)
 \leq\frac{(M-m)\log A}{\log(R/r)},\qquad x\in E.
 \label{eq:nef-radius-comparison}
\end{equation}
For each sufficiently small fixed $r$, the function $x\mapsto S_x(r)$
is plurisubharmonic on its open domain of definition.
\end{lemma}

\begin{proof}
First, $t\mapsto S_x(e^t)$ is nondecreasing and convex.  To see
convexity, fix $0<r_1<r_2\leq R$ and compare the restriction of
$\psi$ to each complex line through $x$ with the harmonic function
on its annulus $r_1<|z-x|<r_2$ that is affine in $\log|z-x|$ and
takes the constant values $S_x(r_1)$ and $S_x(r_2)$ on the boundary
circles.  The subharmonic maximum principle bounds $\psi$ by this
function on each annulus.  On the inner ball the bound follows from
the monotonicity of the interpolation.  Taking the supremum over an
intermediate ball proves convexity.

Apply this convexity at the three logarithmic radii $\log r$,
$\log r+\log A$, and $\log R$.  It gives
\[
 S_x(Ar)-S_x(r)
 \leq\frac{\log A}{\log(R/r)}\bigl(S_x(R)-S_x(r)\bigr),
\]
which implies Equation~\eqref{eq:nef-radius-comparison}.

For fixed $r$, the supremum over the closed translation ball is
upper semicontinuous.  Indeed, for a convergent sequence of centers,
choose maximizing translation vectors and pass to a convergent
subsequence in $\overline B(0,r)$; upper semicontinuity of $\psi$
gives the required upper bound.  It is also a locally bounded
supremum of the plurisubharmonic translations
$x\mapsto\psi(x+\zeta)$, $|\zeta|\leq r$.  Its upper
semicontinuity therefore makes it plurisubharmonic.
\end{proof}

\begin{proof}[Proof of Lemma~\ref{nef:bounded-smoothing}]
\textbf{Local potentials and regularization.}
Choose finitely many coordinate patches $U_i$, smooth real functions
$a_i$ with $\alpha=\ddc a_i$ on neighborhoods of $\overline U_i$,
and open sets
\[
 V_i\Subset W_i\Subset U_i
\]
such that the $V_i$ cover $M$. Such local potentials exist because
$\alpha$ is a smooth closed real $(1,1)$-form. On overlaps,
$a_i-a_j$ is pluriharmonic.

The distributional positivity of $\alpha+\ddc u$ means that
$a_i+u$ has a unique plurisubharmonic representative $\psi_i$.
These representatives are locally bounded: their almost-everywhere
bounds follow from boundedness of $u$ and smoothness of $a_i$, and
the submean inequality and upper semicontinuity extend these bounds
to every point. Moreover,
\[
 \psi_i-\psi_j=a_i-a_j
\]
pointwise on overlaps, since the two plurisubharmonic functions
$\psi_i$ and $\psi_j+a_i-a_j$ agree almost everywhere.
Thus, after replacing $u$ on a set of measure zero, we may write
$\psi_i=a_i+u$ everywhere. The functions $\psi_i$ are bounded above
and below on $U_i$, with bounds uniform over the finite collection.
The positive distances between the compact sets $\overline W_i$
and the boundaries of $U_i$ allow all subsequent regularizations
to be defined on neighborhoods of $\overline W_i$.

In the coordinates of $U_i$, let
\[
 S_{i,b}(x)=\sup_{|y-x|\leq b}\psi_i(y).
\]
Fix a smooth nonnegative kernel of integral one supported in the
unit ball, and let $\rho_{b/4}$ be its rescaling to radius $b/4$.
For all sufficiently small $b>0$, set
\[
 \varphi_{i,b}=S_{i,b}*\rho_{b/4}
\]
near $\overline W_i$.  Lemma~\ref{nef:radius-comparison} and
convolution show that $\varphi_{i,b}$ is smooth and
plurisubharmonic.  If $|y-x|\leq b/4$, then
\[
 B(x,b/2)\subset B(y,b)\subset B(x,2b).
\]
Taking suprema and then averaging proves
\begin{equation}
 S_{i,b/2}(x)\leq\varphi_{i,b}(x)\leq S_{i,2b}(x).
 \label{eq:nef-convolution-bounds}
\end{equation}

\paragraph{Comparison on overlaps.}
We next prove the uniform overlap comparison
\begin{equation}
 \sup_{W_i\cap W_j}
 \left|(\varphi_{i,b}-a_i)-(\varphi_{j,b}-a_j)\right|
 \longrightarrow0\quad\text{as }b\downarrow0.
 \label{eq:nef-overlap-comparison}
\end{equation}
All coordinate changes and their inverses have bounded derivatives
on neighborhoods of the compact overlap sets
$\overline W_i\cap\overline W_j$.  After decreasing the allowed
radius, coordinate balls centered on these overlaps are therefore
comparable by fixed factors.  Thus there are constants $c,C>0$,
independent of the center and of $b$, such that an $i$-coordinate
ball of radius $b$ contains the $j$-coordinate ball of radius $cb$
and is contained in the $j$-coordinate ball of radius $Cb$.
The balls in question lie in $U_i\cap U_j$.

The identity $\psi_i-\psi_j=a_i-a_j$ holds on this intersection,
and the smooth function $a_i-a_j$ varies by $O(b)$ across any such
ball, uniformly in the center.  These facts, together with
Equation~\eqref{eq:nef-convolution-bounds}, give constants
$c_1,c_2,C_2>0$ for which
\[
 S_{j,c_1b}(x)-a_j(x)-C_2b
 \leq\varphi_{i,b}(x)-a_i(x)
 \leq S_{j,c_2b}(x)-a_j(x)+C_2b.
\]
The analogous bounds for $\varphi_{j,b}-a_j$ have radii $b/2$
and $2b$.  Lemma~\ref{nef:radius-comparison}, using a fixed radius
available near each compact overlap, compares all four suprema.
It bounds their differences by $O(1/|\log b|)$.  This proves
Equation~\eqref{eq:nef-overlap-comparison}, in fact with an error
bounded by $O(1/|\log b|)+O(b)$.  All constants can be chosen
uniformly because there are only finitely many pairs of patches.

\paragraph{Gluing the local functions.}
Choose smooth functions $\chi_i$ on neighborhoods of $\overline W_i$
such that
\[
 -3\leq\chi_i\leq0,\qquad
 \chi_i=0\text{ on }V_i,\qquad
 \chi_i=-3\text{ near }\partial W_i.
\]
There is a fixed constant $C_0\geq0$ satisfying
$\ddc\chi_i\geq-C_0k$ for every $i$.  Given $\sigma>0$, choose
$b$ so small that every discrepancy in
Equation~\eqref{eq:nef-overlap-comparison} is less than $\sigma$,
and set on $W_i$
\[
 f_i=\varphi_{i,b}-a_i+\sigma\chi_i.
\]

Here is a precise regularized maximum with the compatibility needed
when the set of available patches changes.  Fix a smooth even
probability density $\eta_\sigma$ on $\R$, supported in
$(-\sigma/4,\sigma/4)$.  For each nonempty finite index set $I$
define
\[
 \mathcal M_I((t_i)_{i\in I})
 =\int_{\R^I}\max_{i\in I}(t_i+s_i)
       \prod_{i\in I}\eta_\sigma(s_i)\,\dd s_i.
\]
This is smooth, convex, nondecreasing in every argument, and
satisfies
\begin{equation}
 \mathcal M_I(t+c\mathbf1)=\mathcal M_I(t)+c.
 \label{eq:nef-max-translation}
\end{equation}
If $t_i$ is more than $\sigma/2$ below the maximum of the other
arguments, then it never attains the shifted maximum in the
integrand.  Its removal leaves exactly $\mathcal M_{I\setminus\{i\}}$,
because the unused probability density integrates to one.

For $x\in M$, put $I(x)=\{i:x\in W_i\}$ and define
\[
 u_\sigma(x)=\mathcal M_{I(x)}((f_i(x))_{i\in I(x)}).
\]
This is globally smooth.  To verify the only issue, fix
$x_0\in\partial W_i$ and choose $j$ with $x_0\in V_j$.
On a neighborhood of $x_0$ within $W_i$, we have
$\chi_i=-3$ and $\chi_j=0$, so the overlap comparison implies
\[
 f_i<f_j-2\sigma.
\]
Thus branch $i$ is inactive throughout this neighborhood, including
under all the allowed shifts.  The exact deletion property just
proved removes it from the formula before crossing $\partial W_i$.
There are finitely many indices, so this argument removes all
indices whose patch boundaries pass through $x_0$, leaving a smooth
formula on a neighborhood of $x_0$.

\paragraph{The curvature estimate.}
Work near any point in a fixed
coordinate patch with local potential $a_j$ for $\alpha$.  Each active branch has
\[
 a_j+f_i=\varphi_{i,b}+(a_j-a_i)+\sigma\chi_i.
\]
The difference $a_j-a_i$ is pluriharmonic, since both functions
are local potentials for $\alpha$. Hence
\begin{equation}
 \ddc(a_j+f_i)\geq-C_0\sigma k.
 \label{eq:nef-branch-curvature}
\end{equation}
By Equation~\eqref{eq:nef-max-translation},
$a_j+u_\sigma$ is the regularized maximum of these branches.
The first derivatives of $\mathcal M_I$ are nonnegative and sum
to one; its Hessian is positive semidefinite.  The complex Hessian
chain rule therefore expresses $\ddc(a_j+u_\sigma)$ as a convex
combination of the branch Hessians plus a nonnegative $(1,1)$-form.
Equation~\eqref{eq:nef-branch-curvature} consequently gives
\[
 \alpha+\ddc u_\sigma=\ddc(a_j+u_\sigma)
 \geq-C_0\sigma k.
\]
Choose $\sigma>0$ with $C_0\sigma\leq\eps$, and rename this
function $u_\eps$. The resulting smooth closed form represents
$[\alpha]$ and has the required lower bound.
\end{proof}
